\section{Training}
\label{sec:training}

\subsection{The five loss terms}
\label{sec:loss_terms}
Let $\mathbf{x}$ be the clean reference, $\mathbf{b}$ the backbone output and $\hat{\mathbf{x}}$ the
wrapper output. The objective of Equation~\eqref{eq:objective} has the following terms.

\textbf{Fidelity.} This term combines the squared error with
the structural similarity index~\cite{wang2004image},
\begin{equation}
\mathcal{L}_{\mathrm{fid}} = \lVert \hat{\mathbf{x}} - \mathbf{x} \rVert_2^2
 + \big(1 - \mathrm{SSIM}(\hat{\mathbf{x}}, \mathbf{x})\big).
\label{eq:loss_fid}
\end{equation}
A dead zone is applied to the peak signal to noise ratio part of this term, so no penalty is charged
while the drop with respect to the backbone stays below $d$ decibels.

\textbf{Clinical improvement.} This term rewards a rise in the clinical measures relative to the
backbone, beginning with the contrast to noise ratio and the tissue contrast index, TCI. The logarithm
of a ratio makes the term scale free,
\begin{equation}
\mathcal{L}_{\mathrm{clin}} = -\log\frac{\mathrm{CNR}(\hat{\mathbf{x}})}{\mathrm{CNR}(\mathbf{b})}
 - \log\frac{\mathrm{TCI}(\hat{\mathbf{x}})}{\mathrm{TCI}(\mathbf{b})}.
\label{eq:loss_clin}
\end{equation}
The same construction is added for the edge preservation index, the equivalent number of looks and
the signal to noise ratio, each with its own weight, so a gain in one measure cannot be bought by a
collapse in another.

\textbf{Property calibration.} The rule layer compares each property score with a threshold, and if the scores drift away from the
thresholds the comparison becomes uninformative, so this term pulls the scores toward the operating
point,
\begin{equation}
\mathcal{L}_{\mathrm{pred}} = \sum_{i=1}^{6}\big(s_i(\hat{\mathbf{x}}) - \tau_i\big)^2 .
\label{eq:loss_pred}
\end{equation}

\textbf{Cooperation.} The wrapper should act where the backbone is unsure and stay quiet where it is
confident. We measure the agreement between the demand map $\mathbf{u}$ and $\boldsymbol{\pi}$, the summed
strength map of the two proposals that the rule layer of Section~\ref{sec:negotiator} reads, with the
Pearson correlation $\rho$, and penalise disagreement,
\begin{equation}
\mathcal{L}_{\mathrm{coop}} = \tfrac{1}{2}\,\mathrm{ReLU}\big(1 - \rho(\mathbf{u}, \boldsymbol{\pi})\big).
\label{eq:loss_coop}
\end{equation}
The Pearson correlation is differentiable everywhere, invariant to an affine change of either map, and
introduces no window size to tune.

\textbf{Edge supervision.} The additive branch is trained to reproduce the gradient that the backbone
lost, measured with the Sobel operator $\nabla$,
\begin{equation}
\mathcal{L}_{\mathrm{edge}} = \big\lVert \nabla\hat{\mathbf{x}} - \nabla\mathbf{x} \big\rVert_2^2 .
\label{eq:loss_edge}
\end{equation}

\subsection{Data and optimisation schedule}
\label{sec:paired_data}
Training uses paired data, a noisy scan and a clean reference averaged over repeated acquisitions,
and the reference enters the fidelity and clinical terms and supervises the edge branch. Testing uses
no clean image at any point, since the six properties read only the noisy input and the backbone
output, and the reference only scores the results. The backbone is trained first, on the training split with the squared error alone, and is then
frozen. The wrapper is then trained in two stages. Only the clinical terms are active at first, so
the correctors learn a useful direction, and the fidelity term switches on at the midpoint, so they
learn how far they may go. Both stages use AdamW with weight decay $10^{-4}$ on $96 \times 96$ crops
in batches of sixteen, eight for DnCNN and SwinIR, for ten epochs in all, with base learning rates of
$2 \times 10^{-4}$ for the corrector networks, $5 \times 10^{-4}$ for the cooperation head and the
strength predictor and $3 \times 10^{-4}$ for the rule layer, and the output layers of the correctors
and the background smoother at two to eight times the corrector rate. The supplementary file restates
the procedure as pseudocode.

Throughout both stages the normalisation statistics of $f$ are held fixed. DnCNN is the one backbone with batch normalisation, and left in training mode 45 of its
buffers move during fitting, so the frozen reference against which every clinical change of that
backbone is measured reads 30.68 dB and a boundary sharpness of 1.031 when the statistics are held
fixed against 31.21 dB and 0.948 when they are not.

\subsection{How the operating point was chosen}
\label{sec:weight_selection}
The terms of Equation~\eqref{eq:objective} carry fixed relative weights. The training script can
also learn them through an uncertainty weighting, and that path is inactive in every run reported
here, which the absence of its parameters from the released checkpoints confirms. Two quantities are selected
per backbone on the validation split, the fidelity dead zone $d$ of Equation~\eqref{eq:loss_fid} and
the rule of Equation~\eqref{eq:gate} that decides where a correction may act. The grid is the two
gate rules crossed with dead zones of 0.3, 0.6 and 0.9 dB, and the tissue bounded gate crossed with
radii of 41 and 61 pixels. It selected the tissue bounded gate at radius 41 and a dead zone of 0.6 dB
for NAFNet, and the percentile gate at a dead zone of 0.9 dB for the other three. Everything else is
held at one value for every backbone and every run. For each backbone every setting is trained at three random seeds
and scored on the validation images. A setting is admissible only if it clears its bounds at its
\emph{worst} seed rather than on average, and among the
admissible settings we take the one with the largest lower confidence bound on its weakest clinical
measure. When no setting clears the first pair of bounds we relax them one step at a time and record
the step that was needed, so a backbone that required a looser condition is visible rather than
hidden. Every source of randomness is seeded, up to the nondeterminism of the convolution kernels, and
the checkpoints chosen on validation are the checkpoints scored on test.

\subsection{Initialisation and the choice of every constant}
\label{sec:constants}
Every constant appears with its value where it acts and has one of four origins. Fixed structural
choices, never tuned, are the sharpness $\kappa = 10$ of
Equation~\eqref{eq:soft_failure}, the steepness 20 of Equation~\eqref{eq:gate}, the window sizes, the
quarter resolution of the gain network, the nominal bound $\pm 0.08$ of the edge branch, the
thresholds $\tau_i$ of Table~\ref{tab:predicates}, taken from the clinical literature and unchanged by
training to four decimals in every run, and the budgets
$B$ and $C$, tied by the identity of Section~\ref{sec:allocation_defect}. Set once on the validation
split while the method was built and never revisited are $\alpha_{\max} = 0.15$, the gain bound
$\pm 0.20$, the gate percentile 0.15, the three tolerances of Section~\ref{sec:safety}, the half
strength background blend and the learning rates. Selected in the grid of
Section~\ref{sec:weight_selection}, on the validation split only, are the dead zone $d$ and the gate
rule. Learned by the objective are $\beta$, $\theta_1$ to $\theta_5$, $\tau_0$, $\beta_0$ and the
scale of the edge branch. Every layer uses its default initialisation with one deliberate exception,
the bias of the brightness channel of the gain network starts at 0.5, so the raw gain begins at
$0.5\,\alpha_{\max}$ in flat tissue rather than on a flat region of the loss. Nothing was ever chosen
on the test split, and Section~\ref{sec:sensitivity} reports how much the results move when the
constants of the rule layer are changed, the direct test of whether any of them does hidden work.


